\documentclass[10pt]{article}
\usepackage[margin=0.7in]{geometry}
\usepackage[T1]{fontenc}
\usepackage{lmodern}
\usepackage{amsmath,amssymb}
\usepackage{enumitem}
\setlist[enumerate]{leftmargin=*, itemsep=7mm}
\setlist[enumerate,2]{label=(\Alph*), leftmargin=2em, itemsep=0.05em, topsep=0.25em}
\pagestyle{plain}

\begin{document}
\begin{center}
  {\large\bfseries Video 4 Quiz: Relations and Functions}\\[0.25em]
  \textit{Choose the best answer for each question.}
\end{center}
\noindent Name:\ \rule{3.2in}{0.4pt}\hfill Date:\ \rule{1.25in}{0.4pt}
\vspace{0.7em}
\begin{enumerate}
  \item When comparing images and preimages of sets, what does the video say a function $f$ can do?
  \begin{enumerate}
    \item $f$ can compress sets, while $f^{-1}$ can expand them.
    \item $f$ can expand sets, while $f^{-1}$ can compress them.
    \item Both $f$ and $f^{-1}$ can only preserve set size.
    \item Neither $f$ nor $f^{-1}$ can map a singleton to multiple elements.
  \end{enumerate}

  \item According to the video, describing functions only by graphs and formulas can lead students to picture what kind of functions?
  \begin{enumerate}
    \item Only ``nice'' functions.
    \item Only functions with finite domains.
    \item Only functions that are one-to-one.
    \item Only functions with real-valued outputs.
  \end{enumerate}

  \item In the birthday example for relations, what does the notation $[a]$ represent?
  \begin{enumerate}
    \item The set of elements equivalent to $a$.
    \item The date on which $a$ was born.
    \item The Cartesian product $A\times A$.
    \item The codomain of the relation.
  \end{enumerate}

  \item What false assertion by Cauchy is mentioned in the video?
  \begin{enumerate}
    \item Every continuous function has an inverse.
    \item A sequence of continuous functions converging everywhere has a continuous limit.
    \item Every function from $X$ to $Y$ is onto.
    \item Every relation is an equivalence relation.
  \end{enumerate}
\end{enumerate}
\end{document}
